\documentclass[11pt]{article}
\usepackage[margin=2.54cm]{geometry}
\usepackage{parskip}
\usepackage{hyperref}
\usepackage{enumitem}

% IEEE Access template fields
\newcommand{\concern}[2]{\noindent\textbf{Reviewer~#1, Concern~\##2}\par}
\newcommand{\concerntext}[1]{\noindent\textit{#1}\par\medskip}
\newcommand{\authorresponse}[1]{\noindent\textbf{Author response:} #1\par\medskip}
\newcommand{\authoraction}[1]{\noindent\textbf{Author action:} We updated the manuscript by #1\par\bigskip\noindent\rule{\linewidth}{0.4pt}\par\medskip}

\begin{document}

% ---------------------------------------------------------------
% Cover letter (IEEE Access template header)
% ---------------------------------------------------------------
\noindent\textbf{Original Manuscript ID:} Access-2025-52745\\
\textbf{Original Article Title:} Topology-Constrained NeuroB-Rep:
Contact-Aware CAD Reconstruction from 3-D Point Clouds

\bigskip
\noindent To: IEEE Access Editor\\
Re: Response to Reviewers

\medskip
\noindent Dear Editor,

Thank you for allowing a resubmission of our manuscript, with an opportunity
to address the reviewers' comments.
We are uploading (a) our point-by-point response to the comments (below),
(b) an updated manuscript with yellow highlighting indicating changes
(as ``Highlighted PDF''), and (c) a clean updated manuscript without
highlights (``Main Manuscript'').

Best regards,\\
J.~W. Kang et al.

\bigskip\noindent\rule{\linewidth}{0.8pt}\bigskip

% ===============================================================
\section*{Reviewer 1}
% ===============================================================

% ---- R1-M1 ----
\concern{1}{1 (Clarity of Contribution vs.\ Point2CAD)}
\concerntext{The distinction between the proposed method and Point2CAD is not
sufficiently clear. A side-by-side comparison table or flowchart would
help readers understand the novelty.}
\authorresponse{We have added Table~2 (Pipeline Comparison), which contrasts
NeuroB-Rep and Point2CAD across seven design dimensions (topology
handling, contact objective, gating schedule, continuity enforcement,
freeform surfaces, export format, and evaluation protocol).
The Introduction (Sec.~I) has been restructured so that the
industrial motivation paragraph precedes the contribution list,
making the positioning against Point2CAD explicit from the outset.}
\authoraction{adding Table~2 in Sec.~IV and restructuring the Introduction
(Sec.~I, paragraphs 3--5).}

% ---- R1-M2 ----
\concern{1}{2 (Evaluation Protocol Consistency)}
\concerntext{The evaluation protocol mixes different scales and metrics.
A rationale for this hybrid approach is needed.}
\authorresponse{We have added a dedicated ``Rationale for Hybrid Evaluation''
paragraph in Sec.~V-B explaining why contact-gap energies are
evaluated in a normalized scale while geometric distances are reported
in millimeters: contact tolerances are scale-invariant design intents,
whereas Chamfer/Hausdorff must be in physical units to be
interpretable by manufacturing engineers. We also clarify the
ICP alignment step and its role in bridging the two scales.}
\authoraction{adding a new paragraph ``Rationale for Hybrid Evaluation''
in Sec.~V-B.}

% ---- R1-M3 ----
\concern{1}{3 (Missing Baseline Comparison)}
\concerntext{Table~3 reports only your method. At least one baseline (e.g.,
Point2CAD) should be included under the same protocol.}
\authorresponse{Table~7 (Geometry Consistency after ICP) now includes a
Point2CAD column. We ran Point2CAD on the same three components
(Impeller, Shaft, Casing) under the identical ICP protocol
with 120 independent repeats per component.
NeuroB-Rep consistently outperforms Point2CAD: geometric accuracy
and surface coincidence errors are 2--3$\times$ lower, while length
errors show 4--8$\times$ improvement.
(Note: The original Table~3 is now Table~7 due to newly added tables.)}
\authoraction{extending Table~7 with a Point2CAD column (Sec.~VI-G).}

% ---- R1-M4 ----
\concern{1}{4 (Failure Cases)}
\concerntext{The paper does not discuss failure modes. Please analyze cases
where the method does not perform well.}
\authorresponse{We have added Sec.~VIII-E ``Failure Modes and Mitigations''
analyzing the cross-boundary misassignment observed in
Figs.~5--6. The primary failure originates in the upstream
PointNet segmentation stage when geometrically similar adjacent
surfaces receive the same label. We describe the downstream
effects (degraded data-term fit, missing contact edges) and
propose three mitigations: adaptive layer thresholds,
boundary-aware resampling, and multi-scale segmentation.}
\authoraction{adding Sec.~VIII-E with failure analysis and three mitigation
strategies.}

% ---- R1-m1 ----
\concern{1}{5 (Abstract Wording)}
\concerntext{``0.01\,mm'' in the abstract should reflect the range across
components.}
\authorresponse{The value has been changed to ``0.01--0.03\,mm'' to reflect
the range across the three components (Impeller, Shaft, Casing).}
\authoraction{updating the abstract and Sec.~VI with the correct range.}

% ---- R1-m2 ----
\concern{1}{6 (Formula Notation)}
\concerntext{The relationship between $d_{ij,k}$ and $s_{ij,k}$ after Eq.~(2)
is unclear.}
\authorresponse{We added a clarifying sentence after Eq.~(2) linking
$d_{ij,k}$ (gap distance) and $s_{ij,k}$ (projected overlap
distance) and explaining their complementary roles in the
contact objective.}
\authoraction{adding a clarifying sentence after Eq.~(2) in Sec.~IV-B.}

% ---- R1-m3 ----
\concern{1}{7 (Figure Reference Order)}
\concerntext{Figures are not referenced in sequential order.}
\authorresponse{All figure references have been checked and corrected to
follow sequential order throughout the manuscript.}
\authoraction{correcting all figure cross-references throughout the manuscript.}

\bigskip\noindent\rule{\linewidth}{0.8pt}\bigskip

% ===============================================================
\section*{Reviewer 2}
% ===============================================================

% ---- R2-C1 ----
\concern{2}{1 (Comparison with Six Cited Methods)}
\concerntext{The paper should compare with Li et al.\ 2019,
ComplexGen, PrimitiveNet, Supervised Fitting, ParSeNet,
and BrepGen, not just cite them.}
\authorresponse{We provide two levels of comparison.
(1)~Table~1 (Feature Comparison) gives a qualitative
comparison across all surveyed methods on seven capability
axes (topology support, freeform surfaces, STEP export, etc.).
(2)~Table~7 provides a full quantitative comparison with
Point2CAD under identical protocol.
A new Sec.~VIII-B (``Scope of Baseline Comparisons'')
explains why the remaining five methods---Li et al.\ 2019~\cite{b27}
(Supervised Fitting), PrimitiveNet~\cite{b28}, ParSeNet~\cite{b19},
ComplexGen~\cite{b1}, and BrepGen~\cite{b4}---cannot be
compared numerically: dataset mismatch (ABC/ShapeNet/synthetic
vs.\ industrial scans), output-format disparity (chain complexes,
latent codes vs.\ STEP), and evaluation-metric incompatibility
(IoU/F-score vs.\ mm-locked Chamfer/Hausdorff + contact-gap).
Point2CAD is the only candidate sharing the same pipeline
architecture, publicly available code, and compatible input format.}
\authoraction{adding Table~1 (Sec.~III), extending Table~7 (Sec.~VI-G),
and adding Sec.~VIII-B explaining the scope of baseline comparisons.}

% ---- R2-C2 ----
\concern{2}{2 (Paper Outline Missing)}
\concerntext{Add a paper outline at the end of the Introduction.}
\authorresponse{A paper-outline paragraph has been added at the end of Sec.~I,
listing all sections and their content.}
\authoraction{adding an outline paragraph at the end of Sec.~I.}

% ---- R2-C3 ----
\concern{2}{3 (Motivation Paragraph Position)}
\concerntext{The industrial motivation should appear before the contribution
list, not after.}
\authorresponse{The motivation paragraph has been moved to precede the
contribution list in Sec.~I.}
\authoraction{reordering the motivation and contribution paragraphs in Sec.~I.}

% ---- R2-C4/C5 ----
\concern{2}{4 \& 5 (Related Work Structure and Depth)}
\concerntext{The Related Work section has a single subsection and lacks
detailed positioning against existing methods.}
\authorresponse{Sec.~III (Related Work) has been restructured into four
subsections: (A)~Point-Cloud Segmentation,
(B)~Primitive Fitting and Assembly, (C)~B-Rep Learning and
Generative Models, and (D)~Hybrid Reconstruction and Point2CAD.
Each subsection discusses the strengths and limitations of prior
methods and explicitly positions our approach.
Table~1 provides a structured feature comparison.
New references have been added: Li et al.\ 2019, PrimitiveNet
(Huang 2021), Maurell et al.\ 2024.}
\authoraction{restructuring Sec.~III into four subsections and adding Table~1.}

% ---- R2-C6 ----
\concern{2}{6 (Face/Edge/Vertex Recovery from Points)}
\concerntext{How are faces, edges, and vertices recovered from the
unstructured point cloud?}
\authorresponse{Sec.~IV-A now includes a dedicated subsubsection
(``From Points to a Patch Graph,'' Sec.~IV-A-2) explaining the kNN
co-label adjacency mechanism: for each point, the $k$~nearest
neighbors are queried; if a neighbor carries a different patch label,
an edge is recorded between the two patches. This yields the face
adjacency graph whose edges define candidate contact boundaries.}
\authoraction{adding subsubsection Sec.~IV-A-2 (``From Points to a Patch Graph'').}

% ---- R2-C7 ----
\concern{2}{7 (Single Subsection in Method)}
\concerntext{Section~III-A contains only one subsection, which violates
standard structuring conventions.}
\authorresponse{A second subsubsection has been added to Sec.~IV-A,
resolving the single-subsection issue.
Sec.~IV-A now contains two subsubsections:
IV-A-1 (Hybrid Fitting with INRs and Primitives) and
IV-A-2 (From Points to a Patch Graph).}
\authoraction{adding subsubsection IV-A-2 to Sec.~IV-A.}

% ---- R2-C8 ----
\concern{2}{8 (Jang et al.\ Explanation)}
\concerntext{The description of Jang et al.\ [12] is too brief.
Expand the explanation of the mesh repair principles.}
\authorresponse{Sec.~II-C now provides a three-step description of Jang et
al.'s method: (i)~detection of interpenetrating triangles via
oriented bounding-box pre-filtering, (ii)~computation of exact
intersection curves between conflicting triangles, and
(iii)~local re-meshing by splitting edges along intersection
curves and re-triangulating the affected patches. The method
operates in a single pass without global re-meshing, making it
suitable as an online diagnostic.}
\authoraction{expanding Sec.~II-C from one sentence to a full paragraph.}

% ---- R2-C9 ----
\concern{2}{9 (Projection Operator Clarity)}
\concerntext{The passage describing the projection operator is hard to follow.}
\authorresponse{The projection operator description has been rewritten for
clarity, with explicit notation for the nearest-point projection
onto each parametric surface and the role of the projection in
computing gap distances.}
\authoraction{rewriting the projection operator passage in Sec.~IV-B.}

% ---- R2-C10 ----
\concern{2}{10 (``Soft Barrier'' Definition)}
\concerntext{Define ``soft barrier'' at first use.}
\authorresponse{The term is now defined in Sec.~II-C (Preliminaries) at its
first occurrence: ``a smooth penalty function
$\log(1+\exp(\tau-s))$ that grows rapidly as the distance~$s$
between projected points on two patches decreases below
threshold~$\tau$.''}
\authoraction{adding the definition of ``soft barrier'' at first use in Sec.~II-C.}

% ---- R2-C11 ----
\concern{2}{11 (Inconsistent Paragraph Labels)}
\concerntext{The ``a:'' style paragraph labels are inconsistent.}
\authorresponse{All paragraph labels have been unified using the standard
\texttt{\textbackslash paragraph\{\}} command throughout the manuscript.}
\authoraction{standardizing all paragraph labels throughout the manuscript.}

% ---- R2-C12 ----
\concern{2}{12 (Fig.~3 Panel Count)}
\concerntext{The caption references panels (a,b,c,d) but only three
panels are visible.}
\authorresponse{This was a caption error. The original figure has always
contained three panels. The caption has been corrected to
reference (a), (b), and (c) only.}
\authoraction{correcting the Fig.~3 caption to reference three panels.}

% ---- R2-C13 ----
\concern{2}{13 (PCA Normalization Justification)}
\concerntext{The PCA canonical transform is used without justification.
An ablation would strengthen the claim.}
\authorresponse{We conducted a PCA ablation experiment (5~seeds, 60~iterations
each) comparing convergence with and without PCA normalization.
Results are reported in Table~6 (Sec.~VI-B). PCA normalization
yields a modest but consistent improvement: Chamfer distance
improves by 0.3\% and Hausdorff by 1.4\%, with lower variance
across seeds. The primary benefit is stabilizing the initial
fit and reducing sensitivity to scan orientation.}
\authoraction{adding Table~6 and Sec.~VI-B (``Effect of PCA Canonical Transform'').}

% ---- R2-C14 ----
\concern{2}{14 (ISO~5459 Datum Explanation)}
\concerntext{The reference to ISO~5459 needs more context.}
\authorresponse{Sec.~V-A now includes an expanded explanation of the
ISO~5459 datum framework: how datum features are established
from physical surfaces, and how this relates to our
evaluation's reference geometry and coordinate alignment via ICP.}
\authoraction{expanding the ISO~5459 explanation in Sec.~V-A.}

% ---- R2-C15 ----
\concern{2}{15 (GPU Implementation Details)}
\concerntext{No information is provided about the GPU implementation.}
\authorresponse{We have added implementation details in Sec.~V-E
(Computational Cost) and Table~4: all optimizations run on a
single NVIDIA RTX~3090 using PyTorch autograd with CUDA-batched
energy evaluations. NeuroB-Rep requires ${\sim}7$\,min / 4.8\,GB
per component vs.\ Point2CAD's ${\sim}4$\,min / 3.2\,GB.
The overhead stems from the contact-gap and continuity terms,
which require pairwise boundary-point queries absent in
Point2CAD's distance-only objective.}
\authoraction{adding Sec.~V-E (Computational Cost) and Table~4.}

% ---- R2-C16 ----
\concern{2}{16 (Artifact URL)}
\concerntext{Provide a public artifact URL for reproducibility.}
\authorresponse{Due to corporate confidentiality and intellectual property
constraints on the core solver algorithms and the raw industrial scan data,
we cannot publicly release the full optimization codebase.
To support independent reproducibility of our reported results, we have
prepared a public artifact at
\url{https://github.com/Juunary/TopoConstrained-NeuroBrep}
containing three items:
(1)~\textbf{Evaluation scripts.}
A standalone Python script (\texttt{eval/evaluate.py}) that computes
Chamfer distance, Hausdorff distance, and Surface Coincidence from any pair
of input meshes, following the protocol of Table~7 (Sec.~V-B).
Dependencies are limited to \texttt{numpy}, \texttt{trimesh}, and
\texttt{scipy}; no proprietary packages are required.
(2)~\textbf{Default configuration file.}
\texttt{configs/default.yaml} lists every hyperparameter reported in
Table~3 (solver iterations, learning rate, gate schedule, loss weights,
etc.), allowing readers to verify all experimental conditions without
access to the solver code itself.
(3)~\textbf{Representative output meshes.}
Final reconstructed PLY files for the Impeller component (NeuroB-Rep
and Point2CAD), together with the reference mesh and a CSV summary of
the Table~7 values, enabling end-to-end metric verification on the
provided sample.}
\authoraction{adding the artifact URL and a structured description of its
contents in Sec.~VII (Reproducibility and Artifact), and releasing
\texttt{configs/default.yaml} to make all reported hyperparameters
independently verifiable.}

% ---- R2-C17 ----
\concern{2}{17 (Table~3 Baseline Column)}
\concerntext{Table~3 should include at least one baseline for comparison.}
\authorresponse{The consistency table (now Table~7 in the revised manuscript)
now includes a Point2CAD column evaluated under the identical ICP
protocol (120 repeats per component). See Concern~\#3 of Reviewer~1 above.}
\authoraction{adding a Point2CAD column to Table~7 (Sec.~VI-G).}

% ---- R2-C18 ----
\concern{2}{18 (``Reproducibility Hygiene'' Wording)}
\concerntext{The term ``reproducibility hygiene'' is informal.}
\authorresponse{The term has been replaced with ``reproducibility protocol''
throughout the manuscript.}
\authoraction{replacing ``reproducibility hygiene'' with ``reproducibility protocol''
throughout Sec.~VII.}

\bigskip\noindent\rule{\linewidth}{0.8pt}\bigskip

% ===============================================================
\section*{Additional Corrections (author-initiated)}
% ===============================================================

During revision, we identified and corrected the following issues
independently of the reviewer comments:

\begin{enumerate}[leftmargin=*, label=\textbf{\arabic*.}]
  \item \textbf{Evaluation repeat count (1{,}000 $\rightarrow$ 120):}
        To ensure a strictly fair, apples-to-apples comparison with the 
        newly added Point2CAD baseline, we re-evaluated both methods under 
        the exact same protocol. Consequently, the evaluation was standardized 
        to 120 independent repeats per component (rather than the 1,000 repeats 
        initially reported for our method alone). The text and Table~7 caption 
        have been updated to reflect this standardized protocol.

  \item \textbf{Fig.~3 caption (4 panels $\rightarrow$ 3 panels).}
        The figure has always contained three panels; the caption
        erroneously referenced four panels (a,b,c,d).
        Corrected to panels (a), (b), (c).

  \item \textbf{Abstract grammar.}
        ``a based on temperature gating mechanism'' corrected to
        ``a temperature-based gating mechanism.''

  \item \textbf{Computational cost reporting (Table~4).}
        Added Table~4 (Sec.~V-E) reporting training time, GPU memory,
        and per-component overhead for both NeuroB-Rep and Point2CAD
        on an NVIDIA RTX~3090.

  \item \textbf{New Acknowledgment section.}
        Added NRF funding acknowledgment before the References,
        per IEEE Access policy.

  \item \textbf{Section and table renumbering.}
        A new Sec.~II (Preliminaries) was inserted and four new tables
        (Tables~1, 2, 4, 6) were added, causing all subsequent section
        and table numbers to shift.
        All cross-references in the revised manuscript have been updated
        accordingly.
\end{enumerate}

\bigskip
\noindent We believe these revisions comprehensively address all reviewer
concerns. We are grateful for the opportunity to strengthen the manuscript
and welcome any further suggestions.

\end{document}
